\documentclass[10pt]{article}
\usepackage[T1]{fontenc}
\usepackage[utf8]{inputenc}
\usepackage[french]{babel}
\usepackage{amsmath,amssymb,mathrsfs,mathpazo}
\usepackage{geometry}
\geometry{a4paper,margin=2cm}
\usepackage{graphicx,array,tabularx,tabularray,colortbl,multirow,multicol,booktabs,enumitem,pifont,ulem,fancybox,diagbox,cancel}
\usepackage{pstricks,pstricks-add,pst-plot,pst-tree,pst-node,pst-text,pst-func,pst-math,pst-bspline,pst-3dplot,pst-eucl,pst-solides3d}
\let\pscalculate\undefined
\usepackage{pst-bezier}
\usepackage{tikz,pgfplots,tkz-tab}
\usetikzlibrary{arrows,arrows.meta,patterns,calc,positioning,trees,decorations.markings,shapes.geometric,shapes.misc,shapes.arrows,intersections,decorations.pathmorphing,decorations.pathreplacing,backgrounds,shadows,babel}
\pgfplotsset{compat=1.18}
\usepgfplotslibrary{fillbetween}
\usepackage[np]{numprint}
\usepackage{listings,makecell,pas-tableur,twoopt}
\let\barre\undefined
\usepackage[pstricks]{tablvar}
\definecolor{darkgreen}{RGB}{0,100,0}
\providecommand{\red}{\color{red}}
\providecommand{\blue}{\color{blue}}
\providecommand{\green}{\color{green}}
\providecommand{\black}{\color{black}}
\providecommand{\white}{\color{white}}
\providecommand{\gray}{\color{gray}}
\providecommand{\cyan}{\color{cyan}}
\providecommand{\magenta}{\color{magenta}}
\providecommand{\yellow}{\color{yellow}}
\providecommand{\violet}{\color{violet}}
\providecommand{\orange}{\color{orange}}
\providecommand{\euro}{\texteuro}
\providecommand{\EUR}{\texteuro}
\providecommand{\ieme}{\textsuperscript{e}}
\providecommand{\degres}{\ensuremath{^\circ}}
\providecommand{\degree}{\ensuremath{^\circ}}
\providecommand{\vv}[1]{\overrightarrow{#1}}
\providecommand{\Coord}[1]{#1}
\providecommand{\decosix}{}
\providecommand{\starredbullet}{\ensuremath{\bullet}}
\pagestyle{empty}
\setlength{\parindent}{0pt}
\usepackage{adjustbox}\DeclareUnicodeCharacter{2236}{:}\DeclareUnicodeCharacter{2190}{\ensuremath{\leftarrow}}
\providecommand{\R}{\mathbb{R}}
\providecommand{\N}{\mathbb{N}}
\providecommand{\Z}{\mathbb{Z}}
\providecommand{\Q}{\mathbb{Q}}
\providecommand{\D}{\mathbb{D}}
\providecommand{\C}{\mathbb{C}}
\providecommand{\e}{{\mathrm{e}}}
\providecommand{\vect}[1]{\overrightarrow{#1}}
\providecommand{\vectt}[1]{\overrightarrow{\mathrm{#1}}}
\providecommand{\vvt}[1]{\overrightarrow{\mathrm{#1}}}
\providecommand{\np}[2][]{{#2~#1}}
\providecommand{\vv}[1]{{\overrightarrow{#1}}}
\renewcommand{\d}{{\mathrm{d}}}
\providecommand{\degres}{{^\circ}}
\providecommand{\degre}{{^\circ}}
\providecommand{\blue}{\color{blue}}
\providecommand{\red}{\color{red}}
\providecommand{\enskip}{\quad}
\providecommand{\llbracket}{\lbrack\!\lbrack}
\providecommand{\rrbracket}{\rbrack\!\rbrack}
\providecommand{\euro}{{\text{€}}}
\providecommand{\no}{{\text{n°}}}
\providecommand{\interval}[2]{{[#1;#2]}}
\providecommand{\intervFO}[2]{{[#1;#2[}}
\providecommand{\dotfill}{\ldots}

\providecommand{\ds}{\displaystyle}
\providecommand{\textding}[1]{#1}
\providecommand{\Oij}{\(\mathrm{O};\vec{\imath},\vec{\jmath}\)}
\providecommand{\Oijk}{\(\mathrm{O};\vec{\imath},\vec{\jmath},\vec{k}\)}
\providecommand{\Ouv}{\(\mathrm{O};\vec{u},\vec{v}\)}
\providecommand{\intervFF}[2]{\ensuremath{[#1;#2]}}
\renewcommand{\euro}{\ifmmode\text{\texteuro}\else\texteuro\fi}
\DeclareUnicodeCharacter{00B0}{\ifmmode^{\circ}\else\textdegree\fi}
\DeclareUnicodeCharacter{20AC}{\ifmmode\text{\texteuro}\else\texteuro\fi}
% Énoncés : Éducation nationale. Transcription et corrigés : APMEP, Denis Vergès et contributeurs. Découpage et classement : CoopMaths, Sébastien Lozano, Éric Elter et contributeurs. Adaptation de la présentation : Maths Explorer.
% CC BY-SA : https://coopmaths.fr/www/about/
% https://coopmaths.fr/alea/static/bac/2026/tex/bac_2026_06_sujet1_antilles_1.tex
\begin{document}
\noindent\textit{APMEP / CoopMaths. CC BY-SA. Présentation adaptée par Maths Explorer.}\par\medskip
\medskip

\subsection*{Partie A}

L'agence de sûreté nucléaire ASNR contrôle les installations des centrales nucléaires afin de garantir leur sécurité et leur conformité. La fréquence des contrôles dépend de l'état du réacteur.

On considère que dans le parc nucléaire français~:
\begin{itemize}[noitemsep]
\item 20~\% des réacteurs sont en arrêt pour maintenance~;
\item 5~\% des réacteurs fonctionnent au ralenti~;
\item le reste des réacteurs du parc nucléaire est en état de fonctionnement normal.
\end{itemize}

On considère également que parmi les réacteurs en arrêt, 85\,\% sont contrôlés et que parmi les réacteurs fonctionnant au ralenti, 60\,\% sont contrôlés.

On choisit au hasard un réacteur dans le parc nucléaire français, et on note~:
\begin{itemize}[noitemsep]
\item $A$ l'évènement \og Le réacteur est en arrêt \fg~;
\item $R$ l'évènement \og Le réacteur fonctionne au ralenti \fg~;
\item $N$ l'évènement \og Le réacteur fonctionne normalement \fg~;
\item $C$ l'évènement \og Le réacteur subit un contrôle \fg.
\end{itemize}

Pour un évènement quelconque $E$, on désigne par $\overline{E}$ son évènement contraire et par $P(E)$ sa probabilité.

On note $x$ la probabilité que le réacteur subisse un contrôle sachant qu'il fonctionne normalement, c'est--à--dire $x = P_N(C)$.

\medskip

\begin{minipage}{0.57\linewidth}
\begin{enumerate}
\item Recopier l'arbre des probabilités ci-contre en complétant les pointillés à l'aide des données de l'énoncé.
\item Calculer la probabilité que le réacteur soit en arrêt et qu'il subisse un contrôle.
\end{enumerate}

On sait de plus que la probabilité qu'un réacteur du parc nucléaire soit contrôlé est égale à $0,35$.

\begin{enumerate}[resume]
\item Montrer que la probabilité $x$ est égale à $0,2$.
\item Le réacteur choisi subit un contrôle. Déterminer la probabilité qu'il fonctionne normalement.
\end{enumerate}
\end{minipage}\hfill
\begin{minipage}{0.4\linewidth}

\pstree[levelsep=2.5,treesep=.7,treemode=R,nodesep=3pt]{\TR{}}
{
\pstree{\TR{$A$}\taput{\ldots}}
	{
	\TR{$C$}\taput{\ldots}
	\TR{$\overline{C}$}\tbput{\ldots}
	}
	\pstree{\TR{$R$}\taput{$\ldots$}}
	{
	\TR{$C$}\taput{\ldots}
	\TR{$\overline{C}$}\tbput{\ldots}
	}
	\pstree{\TR{$N$}\taput{$\ldots$}}
		{
	\TR{$C$}\taput{$x$}
	\TR{$\overline{C}$}\tbput{\ldots}
	}
}


























\end{minipage}


\subsection*{Partie B}

Dans le Sud de la France, on compte 12 réacteurs nucléaires qu'on numérote de 1 à 12.

On sait que lorsque le réacteur fonctionne normalement, sa production journalière est égale à 16~GWh (gigawattheure), lorsque le réacteur fonctionne au ralenti, sa production journalière est réduite à 10~GWh et enfin lorsque le réacteur est à l'arrêt, sa production journalière est nulle.

\medskip

On rappelle que $P(N) = 0,75$, $P(R) = 0,05$ et $P(A) = 0,20$.

\begin{enumerate}
\item On désigne par $X$ la variable aléatoire qui modélise la production journalière en GWh d'un réacteur.

\begin{enumerate}[label=\textbf{\alph*.}]
\item Recopier le tableau ci-dessous et compléter les pointillés donnant la loi de probabilité de la variable aléatoire $X$. On notera $k$ les valeurs prises par la variable aléatoire $X$ et $p_k$ les probabilités égales à $P(X=k)$.

\begin{center}
\begin{tabularx}{0.8\linewidth}{|*{4}{>{\centering \arraybackslash}X|}}\hline
$k$ & $0$ & $\ldots$ & $\ldots$ \\ \hline
$p_k$ & $0,20$ & $\ldots$ & $\ldots$ \\ \hline
\end{tabularx}
\end{center}

\item Montrer que l'espérance $E(X)$ de la variable aléatoire $X$ est égale à $12,5$ et que sa variance $V(X)$ est égale à $40,75$.
\end{enumerate}

\item Pour tout entier naturel $i$ compris entre 1 et 12, on note $X_i$ la variable aléatoire qui modélise la production journalière en GWh du $i$-ème réacteur.

Ainsi, par exemple, $X_2$ désigne la variable aléatoire qui modélise la production journalière en GWh du 2-ième réacteur.

On suppose que les variables $X_1, X_2, \ldots, X_{12}$ suivent la même loi de probabilité que la variable aléatoire $X$, et qu'elles sont indépendantes.

On considère la variable aléatoire $S$ qui donne la production totale journalière des 12 réacteurs. On a ainsi $S = X_1 + X_2 + \cdots + X_{12}$.

On note $E(S)$ l'espérance de la variable aléatoire $S$ et $V(S)$ sa variance.

\begin{enumerate}[label=\textbf{\alph*.}]
\item Montrer que $E(S) = 150$ et $V(S) = 489$.
\item Afin d'éviter le risque de surcharge du réseau et de couvrir l'ensemble des besoins en électricité, la production totale journalière des 12 réacteurs doit être comprise strictement entre 100~GWh et 200~GWh.

Peut-on affirmer que la probabilité d'éviter le risque de surcharge et de couvrir l'ensemble des besoins en électricité est supérieure à $0,80$~? Justifier.
\end{enumerate}
\end{enumerate}

\end{document}
